%% ============================================================
%%  Chapitre 2 — Analyse Technique et Indicateurs Financiers
%% ============================================================
\chapter{Analyse Technique et Indicateurs Financiers}
\label{ch:at}

\section*{Introduction du chapitre}

L'analyse technique est le socle méthodologique sur lequel repose l'ensemble de la
plateforme développée dans ce projet. Avant de présenter le moteur de signaux
(Chapitre~\ref{ch:signaux}) et le moteur de backtest
(Chapitre~\ref{ch:wfo}), il est nécessaire de poser les fondements conceptuels et
mathématiques des vingt familles d'indicateurs constituant l'univers candidat de la
plateforme. Ce chapitre répond à la question~: \textit{sur quoi se fonde techniquement
un signal de trading, et quelles sont les limites de l'analyse technique prise isolément ?}

%% ------------------------------------------------------------------
\section{Fondements de l'analyse technique}
\label{sec:at_fondements}
%% ------------------------------------------------------------------

\subsection{La Dow Theory — les origines}

L'analyse technique trouve ses racines dans les écrits de \textbf{Charles Dow} à la fin
du XIX\textsuperscript{e} siècle. Fondateur du \textit{Wall Street Journal} et
co-créateur du Dow Jones Industrial Average, Dow a formulé plusieurs observations sur
le comportement des marchés~:

\begin{enumerate}
    \item \textbf{Les marchés évoluent en trois tendances imbriquées.}
    La tendance primaire (bull ou bear market) dure de quelques mois à plusieurs années.
    Les corrections secondaires représentent des contremouvements de 33 à 66~\% de la
    tendance primaire. Les fluctuations tertiaires (intraday, hebdomadaires) constituent
    le bruit de court terme.

    \item \textbf{Le volume confirme la tendance.}
    Un mouvement haussier sur volume croissant est fondamentalement plus fiable qu'une
    hausse sur volume décroissant. Ce principe structure directement la catégorie Volume
    de la plateforme.

    \item \textbf{Les tendances persistent jusqu'à un signal de retournement clair.}
    En l'absence de signal contraire, la tendance en cours est présumée maintenue.
    C'est le fondement des stratégies trend-following.
\end{enumerate}

\subsection{Les trois prémisses de l'analyse technique moderne}

Murphy~\cite{murphy1999} formalise les principes de l'AT moderne autour de trois
prémisses fondamentales~:

\begin{description}
    \item[Prémisse 1 — Le marché escompte tout.]
    Le prix de marché reflète instantanément toute l'information disponible~:
    fondamentaux, anticipations des investisseurs, facteurs macro. Cette prémisse
    justifie d'analyser le prix plutôt que les raisons du mouvement.

    \item[Prémisse 2 — Les prix évoluent en tendances.]
    Les marchés ne sont pas purement aléatoires. Ils exhibent des persistances
    (momentum) et des retours à la moyenne documentés empiriquement~\cite{harvey2016}.
    Cette prémisse est le fondement des vingt familles d'indicateurs~: chacune capte
    l'une de ces deux dynamiques.

    \item[Prémisse 3 — L'histoire se répète.]
    Les configurations récurrentes reflètent des comportements humains stables~:
    aversion aux pertes, ancrage cognitif, réactions asymétriques à l'information.
    Cette récurrence rend les indicateurs paramétriques exploitables sur plusieurs cycles.
\end{description}

\subsection{Position par rapport à l'efficience des marchés}

La théorie des marchés efficients (Fama, 1970) postule que les prix reflètent
instantanément toute l'information, rendant l'AT théoriquement inutile. La position
pragmatique retenue dans ce projet est intermédiaire~:

Le MASI est un marché \textit{semi-efficient à faible liquidité}. La littérature
sur les marchés émergents documente des primes de momentum et des anomalies de retour
à la moyenne statistiquement exploitables~\cite{harvey2016}. De plus, l'AT est
utilisée dans ce projet non comme oracle mais comme \textit{filtre systématique}~:
sa valeur réside dans la cohérence et la discipline qu'elle impose, et c'est le
Signal Engine (Chapitre~\ref{ch:signaux}) qui filtre, via une évaluation hors
échantillon, les indicateurs réellement informatifs.

%% ------------------------------------------------------------------
\section{Classification en quatre catégories}
\label{sec:at_classification}
%% ------------------------------------------------------------------

La plateforme organise ses vingt familles en quatre catégories, chacune capturant une
dimension distincte de la dynamique de marché~:

\begin{description}
    \item[\textbf{Tendance (Trend-Following)}]
    Ces indicateurs mesurent la direction et la persistance d'une tendance. Ils
    sur-performent dans les marchés directionnels et sous-performent en marché range,
    car ils sont naturellement en retard sur les retournements.

    \item[\textbf{Momentum}]
    Ces indicateurs mesurent la vitesse et l'accélération du mouvement des prix. Ils
    permettent de capter l'élan d'une tendance ou son essoufflement — complémentant
    les indicateurs de tendance qui ne signalent que la direction.

    \item[\textbf{Oscillation / Mean-Reversion}]
    Ces indicateurs oscillent dans un intervalle borné et détectent les états de
    surachat (prix trop haut par rapport à une norme récente) et de survente (prix
    trop bas). Ils sur-performent dans les marchés sans tendance et permettent d'identifier
    des points de retournement potentiels.

    \item[\textbf{Volume}]
    Le volume est la « deuxième dimension » du marché. Ces indicateurs mesurent la
    conviction derrière un mouvement~: un prix en hausse sur volume décroissant est
    suspect, un prix en hausse sur volume croissant est bien plus fiable. Particulièrement
    pertinents sur le MASI, où les titres hors top~20 présentent une liquidité faible.
\end{description}

La complémentarité de ces quatre catégories est la justification principale du signal
d'ensemble (Chapitre~\ref{ch:signaux}, couche G)~: un consensus positif à travers
plusieurs catégories orthogonales est structurellement plus robuste qu'un signal
isolé.

\figuretodo{ch2_fig1_taxonomie}{Arbre taxonomique des 20 familles d'indicateurs réparties en 4 catégories (Tendance, Momentum, Oscillation, Volume) — 5 familles par catégorie}

%% ------------------------------------------------------------------
\section{Catégorie Tendance — cinq familles}
\label{sec:at_tendance}
%% ------------------------------------------------------------------

Les cinq familles de cette catégorie partagent une logique commune~: elles identifient
si le prix se situe au-dessus ou au-dessous d'une référence calculée sur l'historique
récent, et si les références courtes et longues convergent dans une direction donnée.
Leur signal est ternaire~: $+1$ (haussier), $-1$ (baissier), $0$ (neutre).

\subsection{SMA — Moyenne Mobile Simple}

La SMA est le plus ancien indicateur de tendance. Elle lisse le bruit journalier en
calculant la moyenne arithmétique des $n$ derniers prix de clôture~:

\begin{equation}
    \text{SMA}_n(t) = \frac{1}{n} \sum_{i=0}^{n-1} P_{t-i}
    \label{eq:sma}
\end{equation}

Dans la plateforme, la règle de signal utilisée est \textit{price vs SMA}~: le titre
est haussier si le prix courant est au-dessus de sa SMA, baissier sinon. Cette règle
est paramétrée par la période $n$, dont la valeur optimale varie selon l'horizon~:
courte (5 à 20~jours) pour les signaux court terme, longue (120 à 250~jours) pour
les signaux de fond.

\textbf{Interprétation économique.}
La SMA représente le prix moyen auquel le titre a été échangé sur les $n$ derniers jours.
Si le prix courant est au-dessus, les acheteurs récents sont en gain~: ils ont tendance
à maintenir leur position. Si le prix est en dessous, les acheteurs récents sont en perte
et une pression vendeuse s'exerce naturellement.

\subsection{EMA — Moyenne Mobile Exponentielle}

L'EMA corrige le principal défaut de la SMA~: la pondération uniforme de tous les prix
passés. Elle donne un poids décroissant exponentiellement aux données plus anciennes~:

\begin{equation}
    \text{EMA}_n(t) = \alpha \cdot P_t + (1 - \alpha) \cdot \text{EMA}_n(t-1),
    \quad \alpha = \frac{2}{n+1}
    \label{eq:ema}
\end{equation}

Le facteur de lissage $\alpha$ détermine la réactivité~: une EMA courte (petit $n$)
réagit rapidement aux changements de prix mais génère plus de faux signaux~; une EMA
longue (grand $n$) est plus stable mais plus lente à réagir. Comme pour la SMA,
le signal est~: $+1$ si $P_t > \text{EMA}_n(t)$, $-1$ sinon.

\subsection{EMA Cross — Croisement de deux EMA}

Le croisement de deux EMA de périodes différentes est l'une des stratégies de
tendance les plus utilisées en trading institutionnel. L'idée est de comparer une
EMA rapide ($n_1$) à une EMA lente ($n_2 > n_1$)~:

\begin{equation}
    s_{\text{EMA-cross}}(t) =
    \begin{cases}
        +1 & \text{si } \text{EMA}_{n_1}(t) > \text{EMA}_{n_2}(t) \\
        -1 & \text{si } \text{EMA}_{n_1}(t) < \text{EMA}_{n_2}(t)
    \end{cases}
    \label{eq:ema_cross}
\end{equation}

\textbf{Interprétation économique.}
Quand la moyenne courte passe au-dessus de la longue (\textit{golden cross}), cela
signifie que le momentum récent accélère par rapport à la tendance de fond — signal
d'achat. L'inverse (\textit{death cross}) signale un retournement baissier. En pratique,
ce croisement est plus lent que le signal \textit{price vs EMA} mais génère moins de
faux signaux, car il filtre les fluctuations de court terme.

\figuretodo{ch2_fig2_sma}{Comparaison SMA 20, SMA 50 et EMA 200 sur l'action IAM (2020–2025)~: golden crosses annotés en vert, death crosses en rouge, illustration des faux signaux en période de consolidation}

\subsection{Ichimoku Kinko Hyo}

Le système Ichimoku, développé par Goichi Hosoda dans les années 1930 au Japon, est
le plus complet des indicateurs de tendance. En un seul graphe, il délivre des
informations sur la tendance, le support/résistance dynamique et le momentum.
Il repose sur cinq composantes~:

\begin{align}
    \text{Tenkan-sen}  &= \frac{\max(H_{9}) + \min(L_{9})}{2}
        \quad \text{(ligne de conversion, 9 périodes)} \label{eq:tenkan} \\[4pt]
    \text{Kijun-sen}   &= \frac{\max(H_{26}) + \min(L_{26})}{2}
        \quad \text{(ligne de base, 26 périodes)} \label{eq:kijun} \\[4pt]
    \text{Senkou A}    &= \frac{\text{Tenkan} + \text{Kijun}}{2}
        \quad \text{(décalé de +26 en avant)} \label{eq:senkou_a} \\[4pt]
    \text{Senkou B}    &= \frac{\max(H_{52}) + \min(L_{52})}{2}
        \quad \text{(décalé de +26 en avant)} \label{eq:senkou_b} \\[4pt]
    \text{Chikou}      &= P_t
        \quad \text{(décalé de $-26$ en arrière)} \label{eq:chikou}
\end{align}

Le \textit{Kumo} (nuage) est délimité par Senkou A et Senkou B. Sa position par
rapport au prix donne l'information principale~: prix au-dessus du nuage = tendance
haussière confirmée~; prix dans le nuage = tendance indéterminée~; prix en dessous
du nuage = tendance baissière. L'épaisseur du nuage reflète la force du support/résistance.

\textbf{Signal canonique.} La plateforme extrait~: $+1$ si prix $>$ Kumo (au-dessus
des deux Senkou), $-1$ si prix $<$ Kumo, $0$ dans le nuage.

\subsection{PSAR — Parabolic Stop and Reverse}

Développé par Welles Wilder, le PSAR est un indicateur de suivi de tendance qui
calcule un point stop évoluant de façon parabolique au fil du temps. Sa particularité~:
il est toujours soit au-dessus du prix (trend baissier) soit en dessous (trend haussier).

L'algorithme maintient un point $\text{PSAR}(t)$ et un \textit{facteur d'accélération}
$\alpha(t)$ qui s'incrémente à chaque nouveau sommet (en trend haussier)~:

\begin{align}
    \text{PSAR}(t) &= \text{PSAR}(t-1) + \alpha(t-1) \cdot \bigl(\text{EP}(t-1) - \text{PSAR}(t-1)\bigr)
    \label{eq:psar} \\
    \alpha(t)      &= \min\!\bigl(\alpha(t-1) + \text{step},\; \alpha_{\max}\bigr)
    \label{eq:psar_af}
\end{align}

où EP (\textit{Extreme Point}) est le plus haut atteint pendant le trend haussier en
cours (ou le plus bas en trend baissier), $\text{step}$ est le pas d'incrémentation
(typiquement 0{,}02) et $\alpha_{\max}$ le plafond (typiquement 0{,}2).

\textbf{Interprétation économique.}
Le PSAR représente un stop dynamique~: plus la tendance dure, plus le PSAR se
rapproche du prix, réduisant le risque accepté. Quand le prix touche le PSAR,
le trend est réputé retourné. C'est un des rares indicateurs qui donne simultanément
un signal de direction \textit{et} un niveau de stop.

\textbf{Signal canonique.} $+1$ si PSAR $<$ $P_t$ (PSAR sous le prix = trend haussier),
$-1$ si PSAR $>$ $P_t$.

%% ------------------------------------------------------------------
\section{Catégorie Momentum — cinq familles}
\label{sec:at_momentum}
%% ------------------------------------------------------------------

Ces indicateurs mesurent la vitesse et l'élan du mouvement de prix. Ils répondent à
la question~: \textit{est-ce que le marché accélère ou ralentit dans sa direction
actuelle~?} Leur rôle dans la plateforme est de confirmer ou infirmer les signaux
de tendance~: un signal haussier de tendance confirmé par un momentum fort est bien
plus fiable qu'un signal isolé.

\subsection{MACD — Moving Average Convergence/Divergence}

Développé par Gerald Appel, le MACD mesure la convergence et la divergence entre deux
EMA de périodes différentes. Il est l'un des indicateurs de momentum les plus utilisés
au monde par les traders institutionnels~:

\begin{align}
    \text{MACD}(t)     &= \text{EMA}_{\text{fast}}(t) - \text{EMA}_{\text{slow}}(t)
    \label{eq:macd_line} \\[4pt]
    \text{Signal}(t)   &= \text{EMA}_{\text{signal}}\!\bigl(\text{MACD}(t)\bigr)
    \label{eq:macd_signal} \\[4pt]
    \text{Histogramme}(t) &= \text{MACD}(t) - \text{Signal}(t)
    \label{eq:macd_hist}
\end{align}

Dans la configuration par défaut~: $\text{fast} = 12$, $\text{slow} = 26$,
$\text{signal} = 9$. La plateforme optimise ces trois paramètres lors du WFO.

\textbf{Signal canonique.}
$+1$ si MACD $>$ Signal (l'EMA rapide accélère par rapport à l'EMA lente),
$-1$ si MACD $<$ Signal.

\textbf{La divergence MACD-prix.}
Un phénomène particulièrement pertinent~: si le prix monte mais que le MACD forme
des sommets décroissants, cela révèle un essoufflement du momentum malgré la
progression du prix. Ces divergences sont des signaux avancés de retournement très
utilisés dans l'analyse des valeurs du MASI.

\subsection{ROC — Rate of Change}

Le ROC est la forme la plus directe de mesure du momentum~: il mesure le taux de
variation du prix sur $n$ périodes.

\begin{equation}
    \text{ROC}_n(t) = \frac{P_t - P_{t-n}}{P_{t-n}} \times 100
    \label{eq:roc}
\end{equation}

Un ROC positif et croissant confirme un momentum haussier persistant~; un ROC positif
mais décroissant indique une tendance haussière qui s'essouffle.

\textbf{Signal canonique.}
Croisement de la ligne zéro~: $+1$ si ROC $> 0$, $-1$ si ROC $< 0$.
La plateforme paramétrise $n$ selon l'horizon (5 à 250~jours).

\subsection{TRIX — Triple Exponential Average}

Le TRIX est le taux de variation sur une période d'une EMA triple (EMA appliquée
trois fois successivement)~:

\begin{align}
    \text{EMA1}_n(t) &= \text{EMA}_n(P_t) \label{eq:trix1} \\
    \text{EMA2}_n(t) &= \text{EMA}_n(\text{EMA1}_n(t)) \label{eq:trix2} \\
    \text{EMA3}_n(t) &= \text{EMA}_n(\text{EMA2}_n(t)) \label{eq:trix3} \\
    \text{TRIX}_n(t) &= \frac{\text{EMA3}_n(t) - \text{EMA3}_n(t-1)}{\text{EMA3}_n(t-1)} \times 100
    \label{eq:trix}
\end{align}

Le triple lissage exponentiel élimine la quasi-totalité des cycles court terme, ne
conservant que les tendances significatives. Le TRIX oscille autour de zéro~:
positif = momentum haussier, négatif = baissier.

\textbf{Avantage clé par rapport au ROC.}
Grâce au triple lissage, le TRIX génère beaucoup moins de faux signaux que le ROC
en marché volatile. C'est son principal avantage sur le MASI où certains titres peu
liquides présentent des fluctuations erratiques.

\textbf{Signal canonique.} Croisement de zéro~: $+1$ si TRIX $> 0$, $-1$ si TRIX $< 0$.

\subsection{ADX — Average Directional Index}

L'ADX de Welles Wilder est unique dans cette catégorie~: il mesure la \textit{force}
d'une tendance, quelle que soit sa direction. Il est calculé à partir des Indicateurs
Directionnels $+DI$ et $-DI$~:

\begin{align}
    +DM(t) &= \max(H_t - H_{t-1},\; 0) \quad \text{si } (H_t - H_{t-1}) > (L_{t-1} - L_t) \label{eq:plus_dm} \\
    -DM(t) &= \max(L_{t-1} - L_t,\; 0) \quad \text{si } (L_{t-1} - L_t) > (H_t - H_{t-1}) \label{eq:minus_dm} \\
    +DI_n  &= 100 \cdot \frac{\text{Wilder}_n(+DM)}{\text{ATR}_n} \label{eq:plus_di} \\
    -DI_n  &= 100 \cdot \frac{\text{Wilder}_n(-DM)}{\text{ATR}_n} \label{eq:minus_di} \\
    \text{ADX}_n(t) &= \text{Wilder}_n\!\left(\frac{|+DI_n - (-DI_n)|}{+DI_n + (-DI_n)} \times 100\right)
    \label{eq:adx}
\end{align}

où Wilder$_n$ désigne la moyenne mobile de Wilder (variante de l'EMA avec $\alpha = 1/n$).

\textbf{Interprétation.}
ADX $< 20$~: pas de tendance significative~; ADX entre 20 et 40~: tendance modérée~;
ADX $> 40$~: tendance forte. Le signe de la tendance provient de $+DI$ vs $-DI$~:
$+DI > -DI$ = haussier.

\textbf{Signal canonique.}
La plateforme combine la force (ADX $>$ seuil paramétrable) et la direction~:
$+1$ si ADX $>$ seuil \textbf{et} $+DI > -DI$~;
$-1$ si ADX $>$ seuil \textbf{et} $-DI > +DI$~;
$0$ si ADX $<$ seuil (pas de tendance suffisante).

\subsection{TSI — True Strength Index}

Développé par William Blau, le TSI est un oscillateur de momentum basé sur un double
lissage de la variation de prix~:

\begin{align}
    \Delta P(t)      &= P_t - P_{t-1} \label{eq:tsi_dp} \\
    \text{TSI}(t)    &= 100 \cdot \frac{\text{EMA}_{s}(\text{EMA}_{l}(\Delta P))}
                                        {\text{EMA}_{s}(\text{EMA}_{l}(|\Delta P|))}
    \label{eq:tsi}
\end{align}

où $l$ est la période longue et $s$ la période courte de lissage. Le double lissage
au numérateur capte la direction du momentum, tandis que le dénominateur normalise
par la magnitude absolue des mouvements.

\textbf{Avantage.}
Le TSI oscille dans $[-100, +100]$ et est moins sensible aux gaps de prix que le ROC
ou le MACD, ce qui le rend robuste aux annonces de résultats d'entreprises
(événements fréquents sur le MASI). Son croisement de la ligne zéro est le signal
de changement de momentum.

\textbf{Signal canonique.} $+1$ si TSI $> 0$, $-1$ si TSI $< 0$.

%% ------------------------------------------------------------------
\section{Catégorie Oscillation et Mean-Reversion — cinq familles}
\label{sec:at_oscillation}
%% ------------------------------------------------------------------

Ces indicateurs sont construits pour osciller dans un intervalle borné. Ils ne
mesurent pas la direction de la tendance mais l'état de « sur-extension » du prix
par rapport à un équilibre récent. La logique est inverse de celle des indicateurs
de tendance~: on \textit{achète} les extrêmes de survente et on \textit{vend} les
extrêmes de surachat.

\subsection{RSI — Relative Strength Index}

Introduit par Wilder (1978), le RSI est l'oscillateur le plus utilisé au monde. Il
compare la magnitude des hausses récentes aux baisses récentes sur $n$ périodes~:

\begin{equation}
    \text{RSI}_n(t) = 100 - \frac{100}{1 + RS_n(t)},
    \quad RS_n = \frac{\overline{G}_n}{\overline{L}_n}
    \label{eq:rsi}
\end{equation}

où $\overline{G}_n$ est la moyenne des variations positives et $\overline{L}_n$
la moyenne des variations négatives (en valeur absolue) sur les $n$ dernières périodes.
Le RSI oscille entre 0 et 100.

\textbf{Interprétation économique.}
Un RSI $> 70$ signifie que les hausses ont dominé de façon disproportionnée~: le
titre est en zone de surachat, les vendeurs commencent à l'emporter. Inversement,
RSI $< 30$ indique une survente~: le titre a baissé de façon excessive et un rebond
est probable. Ces seuils sont paramétrables dans la plateforme.

\figuretodo{ch2_fig3_rsi}{RSI 14 périodes sur l'action ATW (2022–2025)~: zones de surachat (>70) et survente (<30) annotées avec les retournements effectifs du cours. Illustration d'une divergence haussière en 2023}

\subsection{Stochastique — Stochastic Oscillator}

Développé par George Lane, l'oscillateur stochastique mesure la position du prix de
clôture par rapport au range (haut/bas) des $n$ dernières périodes. Le raisonnement~:
en tendance haussière, les cours ont tendance à clôturer près des hauts~; en tendance
baissière, près des bas.

\begin{align}
    \%K_n(t) &= \frac{P_t - L_n(t)}{H_n(t) - L_n(t)} \times 100 \label{eq:stoch_k} \\[4pt]
    \%D_n(t) &= \text{SMA}_d\!\bigl(\%K_n\bigr)(t) \label{eq:stoch_d}
\end{align}

où $H_n = \max(P_{t-n+1}, \ldots, P_t)$, $L_n = \min(P_{t-n+1}, \ldots, P_t)$,
et $\%D$ est une moyenne mobile de $\%K$ sur $d$ périodes (ligne signal).

\textbf{Signal canonique.}
$-1$ si $\%K > 80$ (surachat)~; $+1$ si $\%K < 20$ (survente). La plateforme
paramétrise $n$ (période K) et $d$ (période D).

\subsection{CCI — Commodity Channel Index}

Le CCI de Donald Lambert mesure l'écart du prix typique par rapport à sa moyenne
mobile, normalisé par la déviation absolue moyenne~:

\begin{equation}
    \text{CCI}_n(t) = \frac{P^{\text{typ}}_t - \text{SMA}_n(P^{\text{typ}})}{0{,}015 \cdot \text{MAD}_n(t)}
    \label{eq:cci}
\end{equation}

où $P^{\text{typ}}_t = (H_t + L_t + P_t) / 3$ est le prix typique,
$\text{MAD}_n$ la déviation absolue moyenne sur $n$ périodes, et $0{,}015$ un facteur
de calibrage pour que 70--80~\% des valeurs tombent dans $[-100, +100]$.

\textbf{Interprétation.}
CCI $> +100$~: le prix dépasse significativement sa moyenne récente, état de surachat.
CCI $< -100$~: état de survente. La force du CCI réside dans sa réactivité~: en
utilisant la déviation absolue (plus robuste aux outliers que l'écart-type), il
s'adapte automatiquement à la volatilité du titre.

\textbf{Signal canonique.} $-1$ si CCI $> +100$~; $+1$ si CCI $< -100$~; $0$ entre.

\subsection{MFI — Money Flow Index}

Le MFI est souvent qualifié de « RSI pondéré par le volume ». Il intègre le volume
dans l'analyse de l'état de surachat/survente, ce qui lui permet de distinguer une
hausse soutenue par des flux réels d'une hausse spéculative~:

\begin{align}
    P^{\text{typ}}_t &= \frac{H_t + L_t + P_t}{3} \label{eq:mfi_typ} \\[4pt]
    MF_t &= P^{\text{typ}}_t \times V_t \label{eq:mfi_mf} \\[4pt]
    MFR_n &= \frac{\sum_{i : P^{\text{typ}}_i \geq P^{\text{typ}}_{i-1}} MF_i}
                  {\sum_{i : P^{\text{typ}}_i < P^{\text{typ}}_{i-1}} MF_i}
    \quad \text{sur } n \text{ jours} \label{eq:mfi_mfr} \\[4pt]
    \text{MFI}_n(t) &= 100 - \frac{100}{1 + MFR_n}
    \label{eq:mfi}
\end{align}

\textbf{Signal canonique.}
$-1$ si MFI $> 80$ (flux acheteur excessif)~;
$+1$ si MFI $< 20$ (flux vendeur excessif).

\textbf{Pertinence MASI.}
Le MFI est particulièrement utile sur les titres à volume variable du MASI~: une
survente sur volume important (MFI très bas) est un signal bien plus fiable qu'une
survente sur volume anémique, car elle témoigne d'une vraie pression vendeuse
institutionnelle plutôt que d'un simple absence d'acheteurs.

\subsection{UO — Ultimate Oscillator}

Développé par Larry Williams, l'Ultimate Oscillator est conçu pour corriger un défaut
des oscillateurs mono-période~: la sensibilité au choix de la période. Il combine
trois oscillateurs sur trois périodes différentes avec une pondération décroissante~:

\begin{align}
    \text{BP}(t)  &= P_t - \min(L_t, P_{t-1}) \label{eq:uo_bp} \\[4pt]
    \text{TR}(t)  &= \max(H_t, P_{t-1}) - \min(L_t, P_{t-1}) \label{eq:uo_tr} \\[4pt]
    A_{p}         &= \frac{\sum_{i=t-p+1}^{t} \text{BP}(i)}{\sum_{i=t-p+1}^{t} \text{TR}(i)}
    \quad \text{pour } p \in \{p_1, p_2, p_3\} \label{eq:uo_avg} \\[4pt]
    \text{UO}(t)  &= 100 \cdot \frac{4 A_{p_1} + 2 A_{p_2} + A_{p_3}}{4 + 2 + 1}
    \label{eq:uo}
\end{align}

où $p_1 < p_2 < p_3$ (typiquement 7, 14, 28). La pondération 4:2:1 favorise les
périodes courtes (plus réactives) sans ignorer le contexte long terme.

\textbf{Signal canonique.}
Surachat~: UO $> 70$~; survente~: UO $< 30$. La logique de divergence est
particulièrement efficace avec l'UO (signal classique de Williams).

%% ------------------------------------------------------------------
\section{Catégorie Volume — cinq familles}
\label{sec:at_volume}
%% ------------------------------------------------------------------

Les indicateurs de volume répondent à la question fondamentale~: \textit{le mouvement
de prix est-il soutenu par une conviction réelle des participants~?} Un prix qui monte
sur un volume croissant est soutenu par de réels flux acheteurs~; un prix qui monte
sur un volume décroissant trahit une perte de conviction.

Sur le MASI, cette dimension est d'autant plus critique que la liquidité est
concentrée sur une vingtaine de valeurs. Pour les titres peu liquides, un faible volume
peut générer des mouvements de prix violents mais non représentatifs de l'état du marché.

\subsection{OBV — On-Balance Volume}

L'OBV de Joe Granville est le plus ancien indicateur de volume. C'est un compteur
cumulatif qui additionne le volume des séances haussières et soustrait celui des
séances baissières~:

\begin{equation}
    \text{OBV}(t) = \text{OBV}(t-1) +
    \begin{cases}
        +V_t & \text{si } P_t > P_{t-1} \\
        -V_t & \text{si } P_t < P_{t-1} \\
        0    & \text{si } P_t = P_{t-1}
    \end{cases}
    \label{eq:obv}
\end{equation}

\textbf{Interprétation économique.}
Si l'OBV monte tandis que le prix stagne ou baisse légèrement, cela signifie que du
volume s'accumule discrètement en phase d'achat~: un mouvement haussier est en préparation
(accumulation). La plateforme dérive le signal de la pente de l'OBV lissée par une
EMA, ce qui filtre le bruit séance-par-séance.

\textbf{Signal canonique.} $+1$ si pente EMA(OBV) $> 0$~; $-1$ sinon.

\subsection{CMF — Chaikin Money Flow}

Le CMF de Marc Chaikin est une version normalisée et fenêtrée de la ligne
Accumulation/Distribution. Il mesure sur $n$ jours la proportion de flux de capital
qui entre dans le titre par rapport aux sorties~:

\begin{align}
    \text{CLV}(t) &= \frac{(P_t - L_t) - (H_t - P_t)}{H_t - L_t}
        \in [-1, +1] \label{eq:clv} \\[4pt]
    \text{CMF}_n(t) &= \frac{\displaystyle\sum_{i=t-n+1}^{t} \text{CLV}(i) \cdot V_i}
                             {\displaystyle\sum_{i=t-n+1}^{t} V_i}
    \label{eq:cmf}
\end{align}

Le CLV (\textit{Close Location Value}) indique où le prix a clôturé dans le range
de la séance~: $+1$ si au plus haut, $-1$ si au plus bas. Pondéré par le volume et
moyenné sur $n$ jours, il donne une mesure de la pression directionnelle du marché.

\textbf{Signal canonique.} $+1$ si CMF $> 0$ (flux acheteur dominant)~; $-1$ si CMF $< 0$.
Un seuil de $\pm 0{,}05$ est généralement appliqué pour filtrer le bruit.

\subsection{AD — Accumulation/Distribution Line}

La ligne A/D est la version cumulative du CLV~: elle intègre sur toute la durée
disponible la pression directionnelle pondérée par le volume.

\begin{equation}
    \text{AD}(t) = \text{AD}(t-1) + \text{CLV}(t) \times V_t
    \label{eq:ad}
\end{equation}

La valeur absolue de l'A/D n'est pas interprétable~; c'est sa \textit{tendance} qui
importe. La plateforme lisse l'A/D par une EMA et extrait le signal de la pente.

\textbf{Interprétation clé.}
La divergence A/D-prix est l'un des signaux les plus forts de l'analyse par le volume~:
si le prix monte mais que la ligne A/D plafonne ou descend, cela révèle que les
institutionnels sont en train de \textit{distribuer} (vendre) discrètement pendant que
le prix monte sur des flux de détail. Sur le MASI, ce phénomène est observé régulièrement
avant des corrections de titres de grande capitalisation.

\textbf{Signal canonique.} $+1$ si pente EMA(AD) $> 0$~; $-1$ sinon.

\subsection{VWAP — Volume Weighted Average Price}

Le VWAP est la référence de prix institutionnelle universelle. Il représente le prix
moyen auquel le titre a été réellement échangé, pondéré par les volumes~:

\begin{equation}
    \text{VWAP}_n(t) = \frac{\displaystyle\sum_{i=t-n+1}^{t} P^{\text{typ}}_i \cdot V_i}
                             {\displaystyle\sum_{i=t-n+1}^{t} V_i}
    \label{eq:vwap}
\end{equation}

Sur données journalières, la plateforme calcule un VWAP glissant sur $n$ jours. Un prix
au-dessus du VWAP indique que les acheteurs récents ont payé en moyenne plus cher que
le prix d'équilibre~: situation haussière.

\textbf{Pertinence MASI.}
Le VWAP est particulièrement significatif sur les valeurs à forte liquidité (IAM, ATW,
BCP) où les institutionnels exécutent systématiquement leurs ordres autour de ce
niveau. Il est utilisé comme benchmark d'exécution par les desks de trading, ce qui
en fait un niveau de support/résistance dynamique naturel.

\textbf{Signal canonique.} $+1$ si $P_t > \text{VWAP}_n(t) \cdot (1 + \varepsilon)$~;
$-1$ si $P_t < \text{VWAP}_n(t) \cdot (1 - \varepsilon)$, où $\varepsilon$ est un seuil
de déviation paramétrable (typiquement 0{,}5~\% à 2~\%).

\subsection{FI — Force Index (Indice de Force d'Elder)}

Le Force Index d'Alexander Elder mesure la force réelle derrière chaque mouvement de
prix en combinant direction, amplitude et volume~:

\begin{equation}
    \text{FI}(t) = (P_t - P_{t-1}) \times V_t
    \label{eq:fi}
\end{equation}

La plateforme lisse le FI par une EMA de période $n$ pour filtrer le bruit~:
$\text{FI}_n(t) = \text{EMA}_n(\text{FI})$. Un FI positif et croissant indique que
les achats exercent une force croissante~; un FI négatif indique une pression vendeuse.

\textbf{Interprétation.}
Contrairement à l'OBV ou au CMF qui traitent chaque séance de façon binaire (haussière
ou baissière), le FI intègre l'\textit{amplitude} du mouvement~: une hausse de 3~\% sur
fort volume exerce une force bien supérieure à une hausse de 0{,}1~\% sur le même volume.
C'est particulièrement utile sur le MASI pour filtrer les mouvements techniques mineurs
des mouvements à forte conviction.

\textbf{Signal canonique.} $+1$ si EMA(FI) $> 0$~; $-1$ si EMA(FI) $< 0$.

%% ------------------------------------------------------------------
\section{Tableau récapitulatif des vingt familles}
\label{sec:at_tableau}
%% ------------------------------------------------------------------

Le Tableau~\ref{tab:20_familles} récapitule l'ensemble des vingt familles, leur
catégorie, le signal canonique et la période de préchauffage (\textit{warmup})
minimale avant que le signal soit calculable.

\begin{longtable}{p{2.5cm} p{2cm} p{5cm} p{1.5cm}}
\caption{Récapitulatif des 20 familles d'indicateurs de la plateforme}
\label{tab:20_familles} \\
\toprule
\textbf{Famille} & \textbf{Catégorie} & \textbf{Règle de signal ($-1$/$0$/$+1$)} & \textbf{Warmup} \\
\midrule
\endfirsthead
\toprule
\textbf{Famille} & \textbf{Catégorie} & \textbf{Règle de signal ($-1$/$0$/$+1$)} & \textbf{Warmup} \\
\midrule
\endhead
\bottomrule
\endfoot

SMA              & Tendance  & $+1$ si $P_t >$ SMA$_n$                        & $n$ \\
EMA              & Tendance  & $+1$ si $P_t >$ EMA$_n$                        & $n$ \\
EMA Cross        & Tendance  & $+1$ si EMA$_{\text{fast}} >$ EMA$_{\text{slow}}$ & $n_{\text{slow}}$ \\
Ichimoku         & Tendance  & $+1$ si prix $>$ Kumo                          & $n_{\text{senkou\_b}}$ \\
PSAR             & Tendance  & $+1$ si PSAR $< P_t$                           & 2 \\
\midrule
MACD             & Momentum  & $+1$ si MACD $>$ Signal                        & slow + signal \\
ROC              & Momentum  & $+1$ si ROC $> 0$                              & $n$ \\
TRIX             & Momentum  & $+1$ si TRIX $> 0$                             & $3n$ \\
ADX              & Momentum  & $+1$ si ADX $>$ seuil \& $+DI > -DI$          & $2n$ \\
TSI              & Momentum  & $+1$ si TSI $> 0$                              & $l + s$ \\
\midrule
RSI              & Oscillation & $+1$ si RSI $< 30$, $-1$ si RSI $> 70$      & $n$ \\
Stochastique     & Oscillation & $+1$ si \%K $< 20$, $-1$ si \%K $> 80$     & $n_K + n_D$ \\
CCI              & Oscillation & $+1$ si CCI $< -100$, $-1$ si CCI $> +100$ & $n$ \\
MFI              & Oscillation & $+1$ si MFI $< 20$, $-1$ si MFI $> 80$     & $n$ \\
UO               & Oscillation & $+1$ si UO $< 30$, $-1$ si UO $> 70$       & $p_3$ \\
\midrule
OBV              & Volume    & $+1$ si pente EMA(OBV) $> 0$                  & $n$ \\
CMF              & Volume    & $+1$ si CMF $> 0$                              & $n$ \\
AD               & Volume    & $+1$ si pente EMA(A/D) $> 0$                  & $n$ \\
VWAP             & Volume    & $+1$ si $P_t > \text{VWAP}_n \cdot (1+\varepsilon)$ & $n$ \\
FI               & Volume    & $+1$ si EMA(FI) $> 0$                         & $n$ \\

\end{longtable}

%% ------------------------------------------------------------------
\section{Limites de l'AT naïve — pont vers le Signal Engine}
\label{sec:at_limites}
%% ------------------------------------------------------------------

L'analyse technique fournit un vocabulaire riche pour décrire la dynamique des prix.
Cependant, utilisée naïvement — c'est-à-dire par ajustement visuel des paramètres et
lecture subjective de graphes — elle présente des défauts structurels.

\subsection{Paramètres libres et overfitting}

Chaque indicateur possède des paramètres libres. Un analyste peut toujours trouver
\textit{a posteriori} une combinaison de paramètres produisant d'excellents résultats
sur l'historique. C'est ce que Pardo~\cite{pardo2008} appelle la \textit{fitting
exercise}~:

\begin{quote}
\textit{``In-sample backtest is not a backtest — it is a fitting exercise.
The only meaningful measure of a strategy's robustness is its out-of-sample performance.''}
\hfill — Pardo (2008, p.~64)
\end{quote}

\subsection{Tests multiples et faux positifs}

La plateforme teste 20 familles × 30 valeurs de paramètres par horizon × 93 titres
MASI, soit environ 56~000 configurations. Sous un seuil de significativité $\alpha = 5\,\%$,
on s'attend statistiquement à~:

\begin{equation}
    \mathbb{E}\bigl[\text{faux positifs}\bigr] = \alpha \cdot k
    \approx 0{,}05 \times 56\,000 = 2\,800 \text{ indicateurs « significatifs »}
    \label{eq:faux_positifs}
\end{equation}

par pur hasard~\cite{bailey2014}. Ce résultat justifie intégralement l'évaluation
hors-échantillon systématique du Signal Engine.

\subsection{Sensibilité au régime de marché}

Un indicateur de tendance (SMA, PSAR) performe en marché directionnel mais génère
des pertes en marché range. Inversement, un oscillateur (RSI, Stochastique) détecte
bien les retournements mais sur-réagit en tendance forte. Aucun indicateur isolé
n'est robuste à tous les régimes.

\subsection{Synthèse~: de l'AT naïve au Signal Engine}

\begin{table}[H]
\centering
\caption{Problèmes de l'AT naïve et réponses du Signal Engine (Chapitre~\ref{ch:signaux})}
\label{tab:at_vs_signalengine}
\begin{tabularx}{\linewidth}{p{4.5cm} X}
\toprule
\textbf{Problème AT naïve} & \textbf{Réponse Signal Engine} \\
\midrule
Paramètres optimisés in-sample   & Évaluation par sous-périodes indépendantes (couche~B) \\
Tests multiples / faux positifs  & Score multi-composante + floor absolu (couches~C/D) \\
Sensibilité au régime            & Fenêtres multi-régimes + score de stabilité (couche~C) \\
Redondance (corrélation)         & Clustering de Pearson, seuil 0{,}85 (couche~E) \\
Signal unique fragile            & Ensemble pondéré par fiabilité (couche~G) \\
Opacité du raisonnement          & Drill-down complet par famille (dashboard, chapitre~\ref{ch:dashboard}) \\
\bottomrule
\end{tabularx}
\end{table}

\medskip

\noindent Ce chapitre a posé le vocabulaire et la formalisation mathématique des vingt
familles d'indicateurs. Le Chapitre~\ref{ch:signaux} présente la couche de traitement
qui transforme ces familles en un signal d'ensemble robuste et validé hors échantillon.
